\item \textbf{New York Stock Exchange.} Suppose that the percentage returns for a given year for all Stocks listed on the New York Stock Exchange follows a \emph{non}-normal distribution with mean $\mu=11.2$ percent and a standard deviation of $\sigma=8.4$ percent.

\begin{enumerate}
\item  Given the tools we have learned in class, can you find the probability that a \emph{single} stock has an annual return less than 34 percent? Explain why or why not. (Note: you do not need to calculate anything.)\\

\item Consider drawing a random sample of $n=5$ stocks from the population of all stocks and calculating the mean return of the five sampled stocks. Remember that we can think of the \emph{observed} sample mean, $\bar{x}$, as a single draw from a distribution that relates to the random variable $\bar{X}$. We call this distribution \textit{the sampling distribution of the sample mean}.
\begin{enumerate}

\item  What is the shape of the sampling distribution of the sample mean when n=5? Choose the correct answer.%Is it \textsc{normal} or \textsc{non-normal} or is there not enough information to tell? (Multiple Choice) \\
\begin{enumerate}
	\item Normal
	\item Not Normal
	\item Not enough information to tell\\
\end{enumerate}
\item What is the mean of the sampling distribution of the sample mean when n=5? (Report to the nearest two decimal places.)\\
\item What is the standard error of the sampling distribution of the sample mean when n=5? (Report to the nearest two decimal places.)\\
%\item (\textbf{Free Response}) Given the tools we have learned in class, can you find the probability that the mean return of the five sampled stocks is less than 34 percent? Explain why or why not.\\
\item Given the tools we have learned in class, can you find the probability that the mean return of the five sampled stocks is less than 34 percent? Yes or no? Why?\\
\end{enumerate}

\item In general sample size of $n=30$ should be sufficiently large to ensure that $\bar{X}$ follows an approximately Normal distribution. Thus, we are able to calculate probabilities associated with sample means even though we are drawing from a  \emph{non}-normal distribution using the Standard normal table. To answer the following questions you will want to use the following z-score formula:
\begin{equation*}
 z=\frac{\bar{x}-\mu_{\bar{x}}}{\sigma_{\bar{x}}}=\frac{\bar{x}-\mu}{\frac{\sigma}{\sqrt{n}}}.
\end{equation*}
\begin{enumerate}
\item Specify the sampling distribution of the sample mean when $n=30,$ i.e. what are the mean, the standard error and the shape of the sampling distribution?\\
\begin{enumerate}
\item shape: (Multiple Choice)
\begin{itemize}
	\item Normal
	\item Not Normal
	\item Not enough information to tell
\end{itemize}
% \textit{Normal/not Normal/not enough information to tell} \hfill (Multiple Choice.)
\item mean: \underline{\hspace{2cm}} \hfill  (Report to the nearest two decimal places.)
\item standard error: \underline{\hspace{2cm}} \hfill (Report to the nearest two decimal places.) \\
\end{enumerate}
\item Find the probability that the mean return for the 30 sampled
stocks is less than 10 percent. Report the z-score corresponding to 10 percent to 2 decimal places.\\ Report the final probability to 4 decimal places using the standard normal table.\\
\item Find the probability that the mean return for the 30 sampled
stocks is between 10.5 and 14 percent. Report the z-scores corresponding to 10.5 percent and 14 percent to two decimal places. Report the final probability to 4 decimal places using the standard normal table.\\
\item Companies are often concerned with the notion of worst probable case performance of their portfolios. They wish to know the mean return such that we would expect 99\% of samples of 30 stocks to perform better than this mean return. Compute this value and report the answer to two decimal places. \\
\end{enumerate}
\end{enumerate}